\documentclass[12pt]{article}
% Example input for prose_audit.py and outline_check.py: a clean draft (invented content).
\usepackage{amsmath}
\usepackage{booktabs}
\title{Irrigation Canals and Farm Yields}
\author{A. Author}
\begin{document}
\maketitle

\begin{abstract}
Districts with more irrigated land report higher grain yields. The association holds in every
survey year and with district fixed effects.
\end{abstract}

% however --- a comment is not prose and is never audited.
Irrigation is the main public input in dryland farming. Governments spend more on canals than on
any other rural asset. This paper asks how yields move with irrigated area.

\section{Data}
\label{sec:data}

Yields come from an annual farm survey of U.S. districts.\footnote{The survey samples about
\textit{one farm in {forty}}. Weights restore the population totals.} The survey reports output
and planted area for each farm. Dividing output by area gives the yield in tonnes per hectare.

The estimating equation relates the log yield to the irrigated share $s_{dt}$ of district $d$ in
year $t$:
\begin{equation}
\ln y_{dt} = \beta s_{dt} + \alpha_d + \tau_t + \varepsilon_{dt}.
\label{eq:main}
\end{equation}
Here $\alpha_d$ is a district fixed effect, $\tau_t$ a year effect and $\varepsilon_{dt}$ the
error term.

The sample keeps three kinds of districts:
\begin{itemize}
\item districts observed in every survey year;
\item districts with at least ten sampled farms;
\item districts whose borders did not change.
\end{itemize}
Each restriction removes fewer than five percent of the farms.

\section{Results}

A higher irrigated share goes with higher yields. Table~\ref{tab:main} reports the estimates. The
estimate is 0.42 with district and year effects, e.g.\ a rise of ten points in the share goes with
a yield about four percent higher. Rainfall cannot explain the pattern, because the estimate is the
same in wet and dry years.

\begin{table}[htbp]
\centering
\caption{Irrigated share and log yield}
\label{tab:main}
\begin{tabular}{lcc}
\toprule
 & (1) & (2) \\
\midrule
Irrigated share & 0.42 & 0.39 \\
 & (0.08) & --- \\
\bottomrule
\end{tabular}

\smallskip
{\footnotesize \textbf{Notes.} Column (1) includes district and year effects. Standard errors
are clustered by district.}
\end{table}

The estimates come from the survey waves of May 2010 to 2020. The data section describes the
waves in full.

\begin{quote}
However, as one reviewer wrote --- a quotation is not audited.
\end{quote}

\section*{References}
\bibentry Smith, J. 2001. Canals. In \textit{Handbook of Farming}, Chapter 3. However, the
reference list is not audited.
\end{document}
